% ============================================================
% Experiments section for the gravitational clustering paper.
%
% Compile the standalone version with:
%   pdflatex experiments.tex
%
% Or include in a larger document with:
%   % ============================================================
% Experiments section for the gravitational clustering paper.
%
% Compile the standalone version with:
%   pdflatex experiments.tex
%
% Or include in a larger document with:
%   % ============================================================
% Experiments section for the gravitational clustering paper.
%
% Compile the standalone version with:
%   pdflatex experiments.tex
%
% Or include in a larger document with:
%   % ============================================================
% Experiments section for the gravitational clustering paper.
%
% Compile the standalone version with:
%   pdflatex experiments.tex
%
% Or include in a larger document with:
%   \input{experiments}
%
% Required packages (add to the main preamble if using \input):
%   \usepackage{graphicx}
%   \usepackage{amsmath, amssymb}
%   \usepackage{booktabs}
%   \usepackage{subcaption}   % for subfigure / subcaptionbox
%   \usepackage{hyperref}
% ============================================================

\documentclass[11pt, a4paper]{article}

\usepackage{graphicx}
\usepackage{amsmath, amssymb}
\usepackage{booktabs}
\usepackage{subcaption}
\usepackage{hyperref}
\usepackage[margin=2.5cm]{geometry}

% Path to the generated figures (change if the tex and outputs/ live in
% different directories)
\graphicspath{{outputs/}}

\begin{document}

% ============================================================
\section{Experiments}
\label{sec:experiments}
% ============================================================

All experiments were carried out in Python\,3 using NumPy and
scikit-learn.  Unless stated otherwise, data are drawn from a standard
normal distribution in $\mathbb{R}^d$.  The algorithm minimises the
objective
\begin{equation}
  \label{eq:objective}
  F(X) \;=\;
  \sum_{i=1}^n \|x_i - y_i\|^2
  \;+\;
  \frac{2M}{n}
  \sum_{i \neq j}
  \|x_i - x_j\|^2 \,
  e^{-\tfrac{\beta}{2}\|x_i - x_j\|^2},
\end{equation}
via greedy coordinate descent with the analytically derived step size
$\eta = 1/(2 + 4M)$, which is derived from the Lipschitz constant of
$\nabla_{x_i} F$.  The gravitational kernel
$w(t) = e^{-t}(1-t)$, $t = \tfrac{\beta}{2}\|x_i - x_j\|^2$, is
\emph{attractive} for $t < 1$ and \emph{repulsive} for $t > 1$,
creating a natural cluster scale $\sqrt{2/\beta}$.

% ------------------------------------------------------------
\subsection{Experiment~1: Calibration of the Penalty Parameter $M$}
\label{sec:exp1}
% ------------------------------------------------------------

\paragraph{Setup.}
We fix the dimensionality at $d = 2$ and vary the sample size
$n \in \{20, 40, 60, 80\}$, then fix $n = 40$ and vary
$d \in \{2, 3, 4, 5\}$.  For each configuration $(n, d, \beta)$ we
use binary search (8~iterations, starting bracket $[0, 50]$) to find
the minimal penalty $M_{\min}$ such that all points collapse to a
single cluster after 20 gradient steps (collapse criterion: at least
95\% of points within an $\varepsilon = 0.1$ neighbourhood).

\paragraph{Repulsion check.}
If $\beta \cdot 2d > 2.5$ the kernel is entirely repulsive at the
typical inter-point scale and $M_{\min}$ is reported as $\mathrm{NaN}$.

\paragraph{Results.}
Figure~\ref{fig:exp1} shows that $M_{\min}$ grows roughly linearly
with $n$ (left panel) and increases with $d$ (right panel), while
smaller bandwidth $\beta$ consistently requires a larger penalty to
overcome the coarser attraction range.

\begin{figure}[htbp]
  \centering
  \includegraphics[width=\textwidth]{exp1_calibration}
  \caption{%
    \textbf{Experiment~1: Calibration of $M_{\min}$.}
    \emph{Left:} minimal penalty $M_{\min}$ required for full cluster
    collapse as a function of sample size $n$ (fixed $d=2$).
    \emph{Right:} dependence on dimensionality $d$ (fixed $n=40$).
    Three bandwidth values $\beta \in \{0.05, 0.10, 0.20\}$ are shown.
    Points marked as $\mathrm{NaN}$ are omitted (repulsive regime).%
  }
  \label{fig:exp1}
\end{figure}

% ------------------------------------------------------------
\subsection{Experiment~2: Sharp Phase Transition in Two-Cluster
  Separation}
\label{sec:exp2}
% ------------------------------------------------------------

\paragraph{Setup.}
We consider a two-component Gaussian mixture
\begin{equation}
  \label{eq:gmm}
  Y_1 \sim \mathcal{N}(0, I_d), \qquad
  Y_2 \sim \mathcal{N}(r\, e_1, I_d),
\end{equation}
with $n = 100$, $d = 2$, $\beta = 0.1$, $M = 3$.
The separation distance $r$ is swept over
$r \in [1, 7]$ (25 equally spaced values); accuracy is averaged over
three random seeds.  We also visualise the converged point
configurations for $r \in \{2, 4, 6\}$.

\paragraph{Accuracy metric.}
Predicted labels are matched to ground-truth labels via optimal
permutation of the two largest clusters (Hungarian-style over $K=2$).

\paragraph{Results.}
Figure~\ref{fig:exp2} displays a sharp phase transition: below the
critical separation (approximately $r \approx 3.5$--$4.0$) the
algorithm merges both clusters, while above it exact recovery is
achieved.  The scatter plots (top row) confirm that for small $r$ the
two point clouds are dragged into a single mass, whereas for large $r$
each cluster collapses independently to its centroid.

\begin{figure}[htbp]
  \centering
  \includegraphics[width=\textwidth]{exp2_separation}
  \caption{%
    \textbf{Experiment~2: Phase transition in two-cluster separation}
    ($n=100$, $d=2$, $\beta=0.1$, $M=3$).
    \emph{Top row:} converged configurations for separation distances
    $r \in \{2, 4, 6\}$.
    Faint dots — original data; solid diamonds — optimised positions;
    grey lines — displacement trajectories.
    \emph{Bottom:} clustering accuracy as a function of $r$.
    The dashed grey line marks the random-guess baseline (accuracy~$= 0.5$)
    and the dashed green line marks perfect recovery (accuracy~$= 1.0$).%
  }
  \label{fig:exp2}
\end{figure}

% ------------------------------------------------------------
\subsection{Experiment~3: Heterogeneous Clusters and the Need for
  Adaptive $\beta_{ij}$}
\label{sec:exp3}
% ------------------------------------------------------------

\paragraph{Setup.}
We generate a mixture of a \emph{dense} cluster
($n_1 = 50$, $\sigma_1 = 0.2$, centred at the origin) and a
\emph{sparse} cluster ($n_2 = 150$, $\sigma_2 = 1.2$, centred at
$(3.5, 0)$).  Three strategies are compared:
\begin{enumerate}
  \item \textbf{Strategy~A (global small $\beta = 0.03$):}
    the large attraction range absorbs everything into one cluster.
  \item \textbf{Strategy~B (global large $\beta = 0.80$):}
    the fine attraction scale fragments the sparse cluster.
  \item \textbf{Strategy~C (adaptive $\beta_{ij}$ matrix):}
    a symmetric bandwidth matrix is estimated from the $k$-NN
    local scale,
    \begin{equation}
      \label{eq:adaptive_beta}
      \beta_i = \frac{C}{r_k(i)^2}, \qquad
      \beta_{ij} = \sqrt{\beta_i\, \beta_j},
    \end{equation}
    where $r_k(i)$ is the distance to the $k$-th nearest neighbour
    of point $i$, $C = 0.3$, $k = 10$.
    The values are clipped to $[\beta_{\min}, \beta_{\max}] =
    [10^{-4}, 10]$ to prevent numerical issues.
\end{enumerate}
All three strategies use $M = 5$ and 80 gradient steps.

\paragraph{Results.}
Figure~\ref{fig:exp3} shows that only Strategy~C correctly identifies
both clusters; global bandwidth choices fail in qualitatively different
ways.

\begin{figure}[htbp]
  \centering
  \includegraphics[width=\textwidth]{exp3_adaptive}
  \caption{%
    \textbf{Experiment~3: Heterogeneous clusters ($n=200$, $d=2$,
    $M=5$).}
    Grey dots — original data; coloured circles — optimised positions
    coloured by predicted cluster label.
    \emph{Left (Strategy~A, $\beta=0.03$):} large attraction range
    merges both clusters into one.
    \emph{Centre (Strategy~B, $\beta=0.80$):} fine attraction range
    fragments the sparse cluster.
    \emph{Right (Strategy~C, adaptive $\beta_{ij}$):} locally
    adaptive bandwidth correctly recovers both clusters.%
  }
  \label{fig:exp3}
\end{figure}

% ------------------------------------------------------------
\subsection{Experiment~4: Comparison with the Ndaoud~(2020)
  Information-Theoretic Bound}
\label{sec:exp4}
% ------------------------------------------------------------

\paragraph{Theoretical bound.}
Ndaoud~\cite{ndaoud2020} derived the sharp minimax optimal separation
distance for two balanced Gaussian clusters with $\sigma = 1$:
\begin{equation}
  \label{eq:ndaoud}
  r_{\mathrm{opt}}(n, d)
  \;=\;
  2\sqrt{\bar{\Delta}^2}, \qquad
  \bar{\Delta}^2
  = \Bigl(1 + \sqrt{1 + \tfrac{2d}{n\ln n}}\Bigr)\ln n.
\end{equation}

\paragraph{Setup.}
We fix $n = 300$ and vary $d \in \{2, 5, 10, 20, 30, 40\}$.  For each
dimension we set $\beta = 1/d$ (to keep the attraction scale from
shrinking too fast) and $M = 5$.  The empirical minimum separation
$r_{\min}$ is found by binary search (6 steps, bracketed around
$[0.2\,r_{\mathrm{opt}}, 2.5\,r_{\mathrm{opt}}]$), using a strict
exact-recovery criterion: at least 80\% of 10 independent trials must
achieve 100\% classification accuracy.

\paragraph{Results.}
Figure~\ref{fig:exp4} shows that the empirical $r_{\min}$ of the
gravitational algorithm closely tracks the theoretical optimal bound
$r_{\mathrm{opt}}$ across all tested dimensions.  The red-shaded area
between the two curves represents the sub-optimality gap of the
proposed algorithm relative to the information-theoretic limit.

\begin{figure}[htbp]
  \centering
  \includegraphics[width=0.82\textwidth]{exp4_ndaoud_comparison}
  \caption{%
    \textbf{Experiment~4: Comparison with the Ndaoud~(2020)
    information-theoretic lower bound.}
    ($n=300$, $M=5$, $\beta = 1/d$.)
    Dashed black line — theoretical optimal separation $r_{\mathrm{opt}}$
    \eqref{eq:ndaoud}.
    Solid red line — empirical minimal separation $r_{\min}$ of the
    gravitational algorithm.
    Red-shaded region — sub-optimality gap.%
  }
  \label{fig:exp4}
\end{figure}

% ============================================================
% Bibliography (add to the main .bib file when using \input)
% ============================================================
\begin{thebibliography}{9}

\bibitem{ndaoud2020}
  M.~Ndaoud,
  \textit{Sharp optimal recovery in the two component Gaussian mixture model},
  arXiv:2010.04705, 2020.

\end{thebibliography}

\end{document}

%
% Required packages (add to the main preamble if using \input):
%   \usepackage{graphicx}
%   \usepackage{amsmath, amssymb}
%   \usepackage{booktabs}
%   \usepackage{subcaption}   % for subfigure / subcaptionbox
%   \usepackage{hyperref}
% ============================================================

\documentclass[11pt, a4paper]{article}

\usepackage{graphicx}
\usepackage{amsmath, amssymb}
\usepackage{booktabs}
\usepackage{subcaption}
\usepackage{hyperref}
\usepackage[margin=2.5cm]{geometry}

% Path to the generated figures (change if the tex and outputs/ live in
% different directories)
\graphicspath{{outputs/}}

\begin{document}

% ============================================================
\section{Experiments}
\label{sec:experiments}
% ============================================================

All experiments were carried out in Python\,3 using NumPy and
scikit-learn.  Unless stated otherwise, data are drawn from a standard
normal distribution in $\mathbb{R}^d$.  The algorithm minimises the
objective
\begin{equation}
  \label{eq:objective}
  F(X) \;=\;
  \sum_{i=1}^n \|x_i - y_i\|^2
  \;+\;
  \frac{2M}{n}
  \sum_{i \neq j}
  \|x_i - x_j\|^2 \,
  e^{-\tfrac{\beta}{2}\|x_i - x_j\|^2},
\end{equation}
via greedy coordinate descent with the analytically derived step size
$\eta = 1/(2 + 4M)$, which is derived from the Lipschitz constant of
$\nabla_{x_i} F$.  The gravitational kernel
$w(t) = e^{-t}(1-t)$, $t = \tfrac{\beta}{2}\|x_i - x_j\|^2$, is
\emph{attractive} for $t < 1$ and \emph{repulsive} for $t > 1$,
creating a natural cluster scale $\sqrt{2/\beta}$.

% ------------------------------------------------------------
\subsection{Experiment~1: Calibration of the Penalty Parameter $M$}
\label{sec:exp1}
% ------------------------------------------------------------

\paragraph{Setup.}
We fix the dimensionality at $d = 2$ and vary the sample size
$n \in \{20, 40, 60, 80\}$, then fix $n = 40$ and vary
$d \in \{2, 3, 4, 5\}$.  For each configuration $(n, d, \beta)$ we
use binary search (8~iterations, starting bracket $[0, 50]$) to find
the minimal penalty $M_{\min}$ such that all points collapse to a
single cluster after 20 gradient steps (collapse criterion: at least
95\% of points within an $\varepsilon = 0.1$ neighbourhood).

\paragraph{Repulsion check.}
If $\beta \cdot 2d > 2.5$ the kernel is entirely repulsive at the
typical inter-point scale and $M_{\min}$ is reported as $\mathrm{NaN}$.

\paragraph{Results.}
Figure~\ref{fig:exp1} shows that $M_{\min}$ grows roughly linearly
with $n$ (left panel) and increases with $d$ (right panel), while
smaller bandwidth $\beta$ consistently requires a larger penalty to
overcome the coarser attraction range.

\begin{figure}[htbp]
  \centering
  \includegraphics[width=\textwidth]{exp1_calibration}
  \caption{%
    \textbf{Experiment~1: Calibration of $M_{\min}$.}
    \emph{Left:} minimal penalty $M_{\min}$ required for full cluster
    collapse as a function of sample size $n$ (fixed $d=2$).
    \emph{Right:} dependence on dimensionality $d$ (fixed $n=40$).
    Three bandwidth values $\beta \in \{0.05, 0.10, 0.20\}$ are shown.
    Points marked as $\mathrm{NaN}$ are omitted (repulsive regime).%
  }
  \label{fig:exp1}
\end{figure}

% ------------------------------------------------------------
\subsection{Experiment~2: Sharp Phase Transition in Two-Cluster
  Separation}
\label{sec:exp2}
% ------------------------------------------------------------

\paragraph{Setup.}
We consider a two-component Gaussian mixture
\begin{equation}
  \label{eq:gmm}
  Y_1 \sim \mathcal{N}(0, I_d), \qquad
  Y_2 \sim \mathcal{N}(r\, e_1, I_d),
\end{equation}
with $n = 100$, $d = 2$, $\beta = 0.1$, $M = 3$.
The separation distance $r$ is swept over
$r \in [1, 7]$ (25 equally spaced values); accuracy is averaged over
three random seeds.  We also visualise the converged point
configurations for $r \in \{2, 4, 6\}$.

\paragraph{Accuracy metric.}
Predicted labels are matched to ground-truth labels via optimal
permutation of the two largest clusters (Hungarian-style over $K=2$).

\paragraph{Results.}
Figure~\ref{fig:exp2} displays a sharp phase transition: below the
critical separation (approximately $r \approx 3.5$--$4.0$) the
algorithm merges both clusters, while above it exact recovery is
achieved.  The scatter plots (top row) confirm that for small $r$ the
two point clouds are dragged into a single mass, whereas for large $r$
each cluster collapses independently to its centroid.

\begin{figure}[htbp]
  \centering
  \includegraphics[width=\textwidth]{exp2_separation}
  \caption{%
    \textbf{Experiment~2: Phase transition in two-cluster separation}
    ($n=100$, $d=2$, $\beta=0.1$, $M=3$).
    \emph{Top row:} converged configurations for separation distances
    $r \in \{2, 4, 6\}$.
    Faint dots — original data; solid diamonds — optimised positions;
    grey lines — displacement trajectories.
    \emph{Bottom:} clustering accuracy as a function of $r$.
    The dashed grey line marks the random-guess baseline (accuracy~$= 0.5$)
    and the dashed green line marks perfect recovery (accuracy~$= 1.0$).%
  }
  \label{fig:exp2}
\end{figure}

% ------------------------------------------------------------
\subsection{Experiment~3: Heterogeneous Clusters and the Need for
  Adaptive $\beta_{ij}$}
\label{sec:exp3}
% ------------------------------------------------------------

\paragraph{Setup.}
We generate a mixture of a \emph{dense} cluster
($n_1 = 50$, $\sigma_1 = 0.2$, centred at the origin) and a
\emph{sparse} cluster ($n_2 = 150$, $\sigma_2 = 1.2$, centred at
$(3.5, 0)$).  Three strategies are compared:
\begin{enumerate}
  \item \textbf{Strategy~A (global small $\beta = 0.03$):}
    the large attraction range absorbs everything into one cluster.
  \item \textbf{Strategy~B (global large $\beta = 0.80$):}
    the fine attraction scale fragments the sparse cluster.
  \item \textbf{Strategy~C (adaptive $\beta_{ij}$ matrix):}
    a symmetric bandwidth matrix is estimated from the $k$-NN
    local scale,
    \begin{equation}
      \label{eq:adaptive_beta}
      \beta_i = \frac{C}{r_k(i)^2}, \qquad
      \beta_{ij} = \sqrt{\beta_i\, \beta_j},
    \end{equation}
    where $r_k(i)$ is the distance to the $k$-th nearest neighbour
    of point $i$, $C = 0.3$, $k = 10$.
    The values are clipped to $[\beta_{\min}, \beta_{\max}] =
    [10^{-4}, 10]$ to prevent numerical issues.
\end{enumerate}
All three strategies use $M = 5$ and 80 gradient steps.

\paragraph{Results.}
Figure~\ref{fig:exp3} shows that only Strategy~C correctly identifies
both clusters; global bandwidth choices fail in qualitatively different
ways.

\begin{figure}[htbp]
  \centering
  \includegraphics[width=\textwidth]{exp3_adaptive}
  \caption{%
    \textbf{Experiment~3: Heterogeneous clusters ($n=200$, $d=2$,
    $M=5$).}
    Grey dots — original data; coloured circles — optimised positions
    coloured by predicted cluster label.
    \emph{Left (Strategy~A, $\beta=0.03$):} large attraction range
    merges both clusters into one.
    \emph{Centre (Strategy~B, $\beta=0.80$):} fine attraction range
    fragments the sparse cluster.
    \emph{Right (Strategy~C, adaptive $\beta_{ij}$):} locally
    adaptive bandwidth correctly recovers both clusters.%
  }
  \label{fig:exp3}
\end{figure}

% ------------------------------------------------------------
\subsection{Experiment~4: Comparison with the Ndaoud~(2020)
  Information-Theoretic Bound}
\label{sec:exp4}
% ------------------------------------------------------------

\paragraph{Theoretical bound.}
Ndaoud~\cite{ndaoud2020} derived the sharp minimax optimal separation
distance for two balanced Gaussian clusters with $\sigma = 1$:
\begin{equation}
  \label{eq:ndaoud}
  r_{\mathrm{opt}}(n, d)
  \;=\;
  2\sqrt{\bar{\Delta}^2}, \qquad
  \bar{\Delta}^2
  = \Bigl(1 + \sqrt{1 + \tfrac{2d}{n\ln n}}\Bigr)\ln n.
\end{equation}

\paragraph{Setup.}
We fix $n = 300$ and vary $d \in \{2, 5, 10, 20, 30, 40\}$.  For each
dimension we set $\beta = 1/d$ (to keep the attraction scale from
shrinking too fast) and $M = 5$.  The empirical minimum separation
$r_{\min}$ is found by binary search (6 steps, bracketed around
$[0.2\,r_{\mathrm{opt}}, 2.5\,r_{\mathrm{opt}}]$), using a strict
exact-recovery criterion: at least 80\% of 10 independent trials must
achieve 100\% classification accuracy.

\paragraph{Results.}
Figure~\ref{fig:exp4} shows that the empirical $r_{\min}$ of the
gravitational algorithm closely tracks the theoretical optimal bound
$r_{\mathrm{opt}}$ across all tested dimensions.  The red-shaded area
between the two curves represents the sub-optimality gap of the
proposed algorithm relative to the information-theoretic limit.

\begin{figure}[htbp]
  \centering
  \includegraphics[width=0.82\textwidth]{exp4_ndaoud_comparison}
  \caption{%
    \textbf{Experiment~4: Comparison with the Ndaoud~(2020)
    information-theoretic lower bound.}
    ($n=300$, $M=5$, $\beta = 1/d$.)
    Dashed black line — theoretical optimal separation $r_{\mathrm{opt}}$
    \eqref{eq:ndaoud}.
    Solid red line — empirical minimal separation $r_{\min}$ of the
    gravitational algorithm.
    Red-shaded region — sub-optimality gap.%
  }
  \label{fig:exp4}
\end{figure}

% ============================================================
% Bibliography (add to the main .bib file when using \input)
% ============================================================
\begin{thebibliography}{9}

\bibitem{ndaoud2020}
  M.~Ndaoud,
  \textit{Sharp optimal recovery in the two component Gaussian mixture model},
  arXiv:2010.04705, 2020.

\end{thebibliography}

\end{document}

%
% Required packages (add to the main preamble if using \input):
%   \usepackage{graphicx}
%   \usepackage{amsmath, amssymb}
%   \usepackage{booktabs}
%   \usepackage{subcaption}   % for subfigure / subcaptionbox
%   \usepackage{hyperref}
% ============================================================

\documentclass[11pt, a4paper]{article}

\usepackage{graphicx}
\usepackage{amsmath, amssymb}
\usepackage{booktabs}
\usepackage{subcaption}
\usepackage{hyperref}
\usepackage[margin=2.5cm]{geometry}

% Path to the generated figures (change if the tex and outputs/ live in
% different directories)
\graphicspath{{outputs/}}

\begin{document}

% ============================================================
\section{Experiments}
\label{sec:experiments}
% ============================================================

All experiments were carried out in Python\,3 using NumPy and
scikit-learn.  Unless stated otherwise, data are drawn from a standard
normal distribution in $\mathbb{R}^d$.  The algorithm minimises the
objective
\begin{equation}
  \label{eq:objective}
  F(X) \;=\;
  \sum_{i=1}^n \|x_i - y_i\|^2
  \;+\;
  \frac{2M}{n}
  \sum_{i \neq j}
  \|x_i - x_j\|^2 \,
  e^{-\tfrac{\beta}{2}\|x_i - x_j\|^2},
\end{equation}
via greedy coordinate descent with the analytically derived step size
$\eta = 1/(2 + 4M)$, which is derived from the Lipschitz constant of
$\nabla_{x_i} F$.  The gravitational kernel
$w(t) = e^{-t}(1-t)$, $t = \tfrac{\beta}{2}\|x_i - x_j\|^2$, is
\emph{attractive} for $t < 1$ and \emph{repulsive} for $t > 1$,
creating a natural cluster scale $\sqrt{2/\beta}$.

% ------------------------------------------------------------
\subsection{Experiment~1: Calibration of the Penalty Parameter $M$}
\label{sec:exp1}
% ------------------------------------------------------------

\paragraph{Setup.}
We fix the dimensionality at $d = 2$ and vary the sample size
$n \in \{20, 40, 60, 80\}$, then fix $n = 40$ and vary
$d \in \{2, 3, 4, 5\}$.  For each configuration $(n, d, \beta)$ we
use binary search (8~iterations, starting bracket $[0, 50]$) to find
the minimal penalty $M_{\min}$ such that all points collapse to a
single cluster after 20 gradient steps (collapse criterion: at least
95\% of points within an $\varepsilon = 0.1$ neighbourhood).

\paragraph{Repulsion check.}
If $\beta \cdot 2d > 2.5$ the kernel is entirely repulsive at the
typical inter-point scale and $M_{\min}$ is reported as $\mathrm{NaN}$.

\paragraph{Results.}
Figure~\ref{fig:exp1} shows that $M_{\min}$ grows roughly linearly
with $n$ (left panel) and increases with $d$ (right panel), while
smaller bandwidth $\beta$ consistently requires a larger penalty to
overcome the coarser attraction range.

\begin{figure}[htbp]
  \centering
  \includegraphics[width=\textwidth]{exp1_calibration}
  \caption{%
    \textbf{Experiment~1: Calibration of $M_{\min}$.}
    \emph{Left:} minimal penalty $M_{\min}$ required for full cluster
    collapse as a function of sample size $n$ (fixed $d=2$).
    \emph{Right:} dependence on dimensionality $d$ (fixed $n=40$).
    Three bandwidth values $\beta \in \{0.05, 0.10, 0.20\}$ are shown.
    Points marked as $\mathrm{NaN}$ are omitted (repulsive regime).%
  }
  \label{fig:exp1}
\end{figure}

% ------------------------------------------------------------
\subsection{Experiment~2: Sharp Phase Transition in Two-Cluster
  Separation}
\label{sec:exp2}
% ------------------------------------------------------------

\paragraph{Setup.}
We consider a two-component Gaussian mixture
\begin{equation}
  \label{eq:gmm}
  Y_1 \sim \mathcal{N}(0, I_d), \qquad
  Y_2 \sim \mathcal{N}(r\, e_1, I_d),
\end{equation}
with $n = 100$, $d = 2$, $\beta = 0.1$, $M = 3$.
The separation distance $r$ is swept over
$r \in [1, 7]$ (25 equally spaced values); accuracy is averaged over
three random seeds.  We also visualise the converged point
configurations for $r \in \{2, 4, 6\}$.

\paragraph{Accuracy metric.}
Predicted labels are matched to ground-truth labels via optimal
permutation of the two largest clusters (Hungarian-style over $K=2$).

\paragraph{Results.}
Figure~\ref{fig:exp2} displays a sharp phase transition: below the
critical separation (approximately $r \approx 3.5$--$4.0$) the
algorithm merges both clusters, while above it exact recovery is
achieved.  The scatter plots (top row) confirm that for small $r$ the
two point clouds are dragged into a single mass, whereas for large $r$
each cluster collapses independently to its centroid.

\begin{figure}[htbp]
  \centering
  \includegraphics[width=\textwidth]{exp2_separation}
  \caption{%
    \textbf{Experiment~2: Phase transition in two-cluster separation}
    ($n=100$, $d=2$, $\beta=0.1$, $M=3$).
    \emph{Top row:} converged configurations for separation distances
    $r \in \{2, 4, 6\}$.
    Faint dots — original data; solid diamonds — optimised positions;
    grey lines — displacement trajectories.
    \emph{Bottom:} clustering accuracy as a function of $r$.
    The dashed grey line marks the random-guess baseline (accuracy~$= 0.5$)
    and the dashed green line marks perfect recovery (accuracy~$= 1.0$).%
  }
  \label{fig:exp2}
\end{figure}

% ------------------------------------------------------------
\subsection{Experiment~3: Heterogeneous Clusters and the Need for
  Adaptive $\beta_{ij}$}
\label{sec:exp3}
% ------------------------------------------------------------

\paragraph{Setup.}
We generate a mixture of a \emph{dense} cluster
($n_1 = 50$, $\sigma_1 = 0.2$, centred at the origin) and a
\emph{sparse} cluster ($n_2 = 150$, $\sigma_2 = 1.2$, centred at
$(3.5, 0)$).  Three strategies are compared:
\begin{enumerate}
  \item \textbf{Strategy~A (global small $\beta = 0.03$):}
    the large attraction range absorbs everything into one cluster.
  \item \textbf{Strategy~B (global large $\beta = 0.80$):}
    the fine attraction scale fragments the sparse cluster.
  \item \textbf{Strategy~C (adaptive $\beta_{ij}$ matrix):}
    a symmetric bandwidth matrix is estimated from the $k$-NN
    local scale,
    \begin{equation}
      \label{eq:adaptive_beta}
      \beta_i = \frac{C}{r_k(i)^2}, \qquad
      \beta_{ij} = \sqrt{\beta_i\, \beta_j},
    \end{equation}
    where $r_k(i)$ is the distance to the $k$-th nearest neighbour
    of point $i$, $C = 0.3$, $k = 10$.
    The values are clipped to $[\beta_{\min}, \beta_{\max}] =
    [10^{-4}, 10]$ to prevent numerical issues.
\end{enumerate}
All three strategies use $M = 5$ and 80 gradient steps.

\paragraph{Results.}
Figure~\ref{fig:exp3} shows that only Strategy~C correctly identifies
both clusters; global bandwidth choices fail in qualitatively different
ways.

\begin{figure}[htbp]
  \centering
  \includegraphics[width=\textwidth]{exp3_adaptive}
  \caption{%
    \textbf{Experiment~3: Heterogeneous clusters ($n=200$, $d=2$,
    $M=5$).}
    Grey dots — original data; coloured circles — optimised positions
    coloured by predicted cluster label.
    \emph{Left (Strategy~A, $\beta=0.03$):} large attraction range
    merges both clusters into one.
    \emph{Centre (Strategy~B, $\beta=0.80$):} fine attraction range
    fragments the sparse cluster.
    \emph{Right (Strategy~C, adaptive $\beta_{ij}$):} locally
    adaptive bandwidth correctly recovers both clusters.%
  }
  \label{fig:exp3}
\end{figure}

% ------------------------------------------------------------
\subsection{Experiment~4: Comparison with the Ndaoud~(2020)
  Information-Theoretic Bound}
\label{sec:exp4}
% ------------------------------------------------------------

\paragraph{Theoretical bound.}
Ndaoud~\cite{ndaoud2020} derived the sharp minimax optimal separation
distance for two balanced Gaussian clusters with $\sigma = 1$:
\begin{equation}
  \label{eq:ndaoud}
  r_{\mathrm{opt}}(n, d)
  \;=\;
  2\sqrt{\bar{\Delta}^2}, \qquad
  \bar{\Delta}^2
  = \Bigl(1 + \sqrt{1 + \tfrac{2d}{n\ln n}}\Bigr)\ln n.
\end{equation}

\paragraph{Setup.}
We fix $n = 300$ and vary $d \in \{2, 5, 10, 20, 30, 40\}$.  For each
dimension we set $\beta = 1/d$ (to keep the attraction scale from
shrinking too fast) and $M = 5$.  The empirical minimum separation
$r_{\min}$ is found by binary search (6 steps, bracketed around
$[0.2\,r_{\mathrm{opt}}, 2.5\,r_{\mathrm{opt}}]$), using a strict
exact-recovery criterion: at least 80\% of 10 independent trials must
achieve 100\% classification accuracy.

\paragraph{Results.}
Figure~\ref{fig:exp4} shows that the empirical $r_{\min}$ of the
gravitational algorithm closely tracks the theoretical optimal bound
$r_{\mathrm{opt}}$ across all tested dimensions.  The red-shaded area
between the two curves represents the sub-optimality gap of the
proposed algorithm relative to the information-theoretic limit.

\begin{figure}[htbp]
  \centering
  \includegraphics[width=0.82\textwidth]{exp4_ndaoud_comparison}
  \caption{%
    \textbf{Experiment~4: Comparison with the Ndaoud~(2020)
    information-theoretic lower bound.}
    ($n=300$, $M=5$, $\beta = 1/d$.)
    Dashed black line — theoretical optimal separation $r_{\mathrm{opt}}$
    \eqref{eq:ndaoud}.
    Solid red line — empirical minimal separation $r_{\min}$ of the
    gravitational algorithm.
    Red-shaded region — sub-optimality gap.%
  }
  \label{fig:exp4}
\end{figure}

% ============================================================
% Bibliography (add to the main .bib file when using \input)
% ============================================================
\begin{thebibliography}{9}

\bibitem{ndaoud2020}
  M.~Ndaoud,
  \textit{Sharp optimal recovery in the two component Gaussian mixture model},
  arXiv:2010.04705, 2020.

\end{thebibliography}

\end{document}

%
% Required packages (add to the main preamble if using \input):
%   \usepackage{graphicx}
%   \usepackage{amsmath, amssymb}
%   \usepackage{booktabs}
%   \usepackage{subcaption}   % for subfigure / subcaptionbox
%   \usepackage{hyperref}
% ============================================================

\documentclass[11pt, a4paper]{article}

\usepackage{graphicx}
\usepackage{amsmath, amssymb}
\usepackage{booktabs}
\usepackage{subcaption}
\usepackage{hyperref}
\usepackage[margin=2.5cm]{geometry}

% Path to the generated figures (change if the tex and outputs/ live in
% different directories)
\graphicspath{{outputs/}}

\begin{document}

% ============================================================
\section{Experiments}
\label{sec:experiments}
% ============================================================

All experiments were carried out in Python\,3 using NumPy and
scikit-learn.  Unless stated otherwise, data are drawn from a standard
normal distribution in $\mathbb{R}^d$.  The algorithm minimises the
objective
\begin{equation}
  \label{eq:objective}
  F(X) \;=\;
  \sum_{i=1}^n \|x_i - y_i\|^2
  \;+\;
  \frac{2M}{n}
  \sum_{i \neq j}
  \|x_i - x_j\|^2 \,
  e^{-\tfrac{\beta}{2}\|x_i - x_j\|^2},
\end{equation}
via greedy coordinate descent with the analytically derived step size
$\eta = 1/(2 + 4M)$, which is derived from the Lipschitz constant of
$\nabla_{x_i} F$.  The gravitational kernel
$w(t) = e^{-t}(1-t)$, $t = \tfrac{\beta}{2}\|x_i - x_j\|^2$, is
\emph{attractive} for $t < 1$ and \emph{repulsive} for $t > 1$,
creating a natural cluster scale $\sqrt{2/\beta}$.

% ------------------------------------------------------------
\subsection{Experiment~1: Calibration of the Penalty Parameter $M$}
\label{sec:exp1}
% ------------------------------------------------------------

\paragraph{Setup.}
We fix the dimensionality at $d = 2$ and vary the sample size
$n \in \{20, 40, 60, 80\}$, then fix $n = 40$ and vary
$d \in \{2, 3, 4, 5\}$.  For each configuration $(n, d, \beta)$ we
use binary search (8~iterations, starting bracket $[0, 50]$) to find
the minimal penalty $M_{\min}$ such that all points collapse to a
single cluster after 20 gradient steps (collapse criterion: at least
95\% of points within an $\varepsilon = 0.1$ neighbourhood).

\paragraph{Repulsion check.}
If $\beta \cdot 2d > 2.5$ the kernel is entirely repulsive at the
typical inter-point scale and $M_{\min}$ is reported as $\mathrm{NaN}$.

\paragraph{Results.}
Figure~\ref{fig:exp1} shows that $M_{\min}$ grows roughly linearly
with $n$ (left panel) and increases with $d$ (right panel), while
smaller bandwidth $\beta$ consistently requires a larger penalty to
overcome the coarser attraction range.

\begin{figure}[htbp]
  \centering
  \includegraphics[width=\textwidth]{exp1_calibration}
  \caption{%
    \textbf{Experiment~1: Calibration of $M_{\min}$.}
    \emph{Left:} minimal penalty $M_{\min}$ required for full cluster
    collapse as a function of sample size $n$ (fixed $d=2$).
    \emph{Right:} dependence on dimensionality $d$ (fixed $n=40$).
    Three bandwidth values $\beta \in \{0.05, 0.10, 0.20\}$ are shown.
    Points marked as $\mathrm{NaN}$ are omitted (repulsive regime).%
  }
  \label{fig:exp1}
\end{figure}

% ------------------------------------------------------------
\subsection{Experiment~2: Sharp Phase Transition in Two-Cluster
  Separation}
\label{sec:exp2}
% ------------------------------------------------------------

\paragraph{Setup.}
We consider a two-component Gaussian mixture
\begin{equation}
  \label{eq:gmm}
  Y_1 \sim \mathcal{N}(0, I_d), \qquad
  Y_2 \sim \mathcal{N}(r\, e_1, I_d),
\end{equation}
with $n = 100$, $d = 2$, $\beta = 0.1$, $M = 3$.
The separation distance $r$ is swept over
$r \in [1, 7]$ (25 equally spaced values); accuracy is averaged over
three random seeds.  We also visualise the converged point
configurations for $r \in \{2, 4, 6\}$.

\paragraph{Accuracy metric.}
Predicted labels are matched to ground-truth labels via optimal
permutation of the two largest clusters (Hungarian-style over $K=2$).

\paragraph{Results.}
Figure~\ref{fig:exp2} displays a sharp phase transition: below the
critical separation (approximately $r \approx 3.5$--$4.0$) the
algorithm merges both clusters, while above it exact recovery is
achieved.  The scatter plots (top row) confirm that for small $r$ the
two point clouds are dragged into a single mass, whereas for large $r$
each cluster collapses independently to its centroid.

\begin{figure}[htbp]
  \centering
  \includegraphics[width=\textwidth]{exp2_separation}
  \caption{%
    \textbf{Experiment~2: Phase transition in two-cluster separation}
    ($n=100$, $d=2$, $\beta=0.1$, $M=3$).
    \emph{Top row:} converged configurations for separation distances
    $r \in \{2, 4, 6\}$.
    Faint dots — original data; solid diamonds — optimised positions;
    grey lines — displacement trajectories.
    \emph{Bottom:} clustering accuracy as a function of $r$.
    The dashed grey line marks the random-guess baseline (accuracy~$= 0.5$)
    and the dashed green line marks perfect recovery (accuracy~$= 1.0$).%
  }
  \label{fig:exp2}
\end{figure}

% ------------------------------------------------------------
\subsection{Experiment~3: Heterogeneous Clusters and the Need for
  Adaptive $\beta_{ij}$}
\label{sec:exp3}
% ------------------------------------------------------------

\paragraph{Setup.}
We generate a mixture of a \emph{dense} cluster
($n_1 = 50$, $\sigma_1 = 0.2$, centred at the origin) and a
\emph{sparse} cluster ($n_2 = 150$, $\sigma_2 = 1.2$, centred at
$(3.5, 0)$).  Three strategies are compared:
\begin{enumerate}
  \item \textbf{Strategy~A (global small $\beta = 0.03$):}
    the large attraction range absorbs everything into one cluster.
  \item \textbf{Strategy~B (global large $\beta = 0.80$):}
    the fine attraction scale fragments the sparse cluster.
  \item \textbf{Strategy~C (adaptive $\beta_{ij}$ matrix):}
    a symmetric bandwidth matrix is estimated from the $k$-NN
    local scale,
    \begin{equation}
      \label{eq:adaptive_beta}
      \beta_i = \frac{C}{r_k(i)^2}, \qquad
      \beta_{ij} = \sqrt{\beta_i\, \beta_j},
    \end{equation}
    where $r_k(i)$ is the distance to the $k$-th nearest neighbour
    of point $i$, $C = 0.3$, $k = 10$.
    The values are clipped to $[\beta_{\min}, \beta_{\max}] =
    [10^{-4}, 10]$ to prevent numerical issues.
\end{enumerate}
All three strategies use $M = 5$ and 80 gradient steps.

\paragraph{Results.}
Figure~\ref{fig:exp3} shows that only Strategy~C correctly identifies
both clusters; global bandwidth choices fail in qualitatively different
ways.

\begin{figure}[htbp]
  \centering
  \includegraphics[width=\textwidth]{exp3_adaptive}
  \caption{%
    \textbf{Experiment~3: Heterogeneous clusters ($n=200$, $d=2$,
    $M=5$).}
    Grey dots — original data; coloured circles — optimised positions
    coloured by predicted cluster label.
    \emph{Left (Strategy~A, $\beta=0.03$):} large attraction range
    merges both clusters into one.
    \emph{Centre (Strategy~B, $\beta=0.80$):} fine attraction range
    fragments the sparse cluster.
    \emph{Right (Strategy~C, adaptive $\beta_{ij}$):} locally
    adaptive bandwidth correctly recovers both clusters.%
  }
  \label{fig:exp3}
\end{figure}

% ------------------------------------------------------------
\subsection{Experiment~4: Comparison with the Ndaoud~(2020)
  Information-Theoretic Bound}
\label{sec:exp4}
% ------------------------------------------------------------

\paragraph{Theoretical bound.}
Ndaoud~\cite{ndaoud2020} derived the sharp minimax optimal separation
distance for two balanced Gaussian clusters with $\sigma = 1$:
\begin{equation}
  \label{eq:ndaoud}
  r_{\mathrm{opt}}(n, d)
  \;=\;
  2\sqrt{\bar{\Delta}^2}, \qquad
  \bar{\Delta}^2
  = \Bigl(1 + \sqrt{1 + \tfrac{2d}{n\ln n}}\Bigr)\ln n.
\end{equation}

\paragraph{Setup.}
We fix $n = 300$ and vary $d \in \{2, 5, 10, 20, 30, 40\}$.  For each
dimension we set $\beta = 1/d$ (to keep the attraction scale from
shrinking too fast) and $M = 5$.  The empirical minimum separation
$r_{\min}$ is found by binary search (6 steps, bracketed around
$[0.2\,r_{\mathrm{opt}}, 2.5\,r_{\mathrm{opt}}]$), using a strict
exact-recovery criterion: at least 80\% of 10 independent trials must
achieve 100\% classification accuracy.

\paragraph{Results.}
Figure~\ref{fig:exp4} shows that the empirical $r_{\min}$ of the
gravitational algorithm closely tracks the theoretical optimal bound
$r_{\mathrm{opt}}$ across all tested dimensions.  The red-shaded area
between the two curves represents the sub-optimality gap of the
proposed algorithm relative to the information-theoretic limit.

\begin{figure}[htbp]
  \centering
  \includegraphics[width=0.82\textwidth]{exp4_ndaoud_comparison}
  \caption{%
    \textbf{Experiment~4: Comparison with the Ndaoud~(2020)
    information-theoretic lower bound.}
    ($n=300$, $M=5$, $\beta = 1/d$.)
    Dashed black line — theoretical optimal separation $r_{\mathrm{opt}}$
    \eqref{eq:ndaoud}.
    Solid red line — empirical minimal separation $r_{\min}$ of the
    gravitational algorithm.
    Red-shaded region — sub-optimality gap.%
  }
  \label{fig:exp4}
\end{figure}

% ============================================================
% Bibliography (add to the main .bib file when using \input)
% ============================================================
\begin{thebibliography}{9}

\bibitem{ndaoud2020}
  M.~Ndaoud,
  \textit{Sharp optimal recovery in the two component Gaussian mixture model},
  arXiv:2010.04705, 2020.

\end{thebibliography}

\end{document}
